% =============================================================================
% [Project title] : paper draft
%
% This is a starter LaTeX template that mirrors the structure of doc/paper.md.
% Each section here corresponds to a § in paper.md; the rules for what goes in
% each section live in doc/CLAUDE.md.
%
% Compile on Overleaf: upload the whole `paper/` folder as a new Overleaf
% project. Compiler: pdfLaTeX. Main document: main.tex.
% =============================================================================
\documentclass[11pt, letterpaper]{article}

\usepackage[margin=1in]{geometry}
\usepackage{amsmath, amssymb}
\usepackage{graphicx}
\usepackage{booktabs}
\usepackage{xcolor}
\usepackage[colorlinks=true, linkcolor=blue!50!black, citecolor=blue!50!black,
            urlcolor=blue!50!black]{hyperref}
\usepackage{enumitem}
\usepackage{titlesec}
\usepackage{microtype}
\usepackage{caption}
\usepackage{authblk}
\usepackage{natbib}
\usepackage{lmodern}   % scalable fonts (required by microtype font expansion)
\usepackage[textwidth=2.5cm, textsize=footnotesize]{todonotes}
\usepackage{placeins}  % \FloatBarrier : keep figures from floating past section boundaries

\titleformat{\section}{\large\bfseries}{\thesection.}{0.5em}{}
\titleformat{\subsection}{\normalsize\bfseries}{\thesubsection}{0.5em}{}

% =============================================================================
% Coauthor comment macros (for review markup; visible in margin)
%
% Usage in the body:
%   \tynote{question or note from the primary author}
%   \robinnote{advisor's comment}
%   \coauthornote{Initials}{generic coauthor's comment}
%
% Inline form (no margin) : when sentence-level annotation is more natural:
%   \inlinecomment{Init}{sentence-level remark}
%
% To strip all comments for a clean draft, replace the todonotes import with
%   \usepackage[disable]{todonotes}
% =============================================================================
\newcommand{\tynote}[1]{\todo[color=blue!25, inline=false]{\textbf{TY:} #1}}
\newcommand{\robinnote}[1]{\todo[color=red!25, inline=false]{\textbf{Robin:} #1}}
\newcommand{\coauthornote}[2]{\todo[color=green!25, inline=false]{\textbf{#1:} #2}}
\newcommand{\inlinecomment}[2]{\textcolor{purple}{\textbf{[#1: }#2\textbf{]}}}

\title{{[Project title]}\\
  {[Subtitle: the one-sentence thesis]}}
\author{[Author]}
\affil{[Affiliation]}
\date{\today}

\begin{document}
\maketitle

\begin{abstract}
[Prior work claim in one sentence.] [Our refined measurement in one
sentence.] [Headline numbers.] [Mechanism story in one sentence.]
[Implication for the field.]

This abstract mirrors the structure of \texttt{doc/paper.md}'s
\textsc{Thesis} subsection. Replace bracketed placeholders with the
project specifics.
\end{abstract}

% =============================================================================
\section{Introduction}
\label{sec:intro}

% House rule (write-paper references/style.md): the introduction has NO subsections —
% flowing prose ending in a numbered contributions list. The five-move order lives in
% doc/example_papers/style_extraction.md §3a.

[Move 1, problem with one worked instance: plain-language setup, what
prior work claimed, the operational question this paper answers. 3--6
sentences a non-specialist can follow.]

[Move 2, the gap: why the obvious experiment or prior fixes do not
settle it. Most projects have a discriminator story; set up the two
theories here.]

\begin{table}[t]
\centering
\small
\caption{Two theories the data are consistent with.}
\label{tab:theories}
\begin{tabular}{p{0.18\textwidth}p{0.36\textwidth}p{0.36\textwidth}}
\toprule
\textbf{Theory} & \textbf{What it predicts} & \textbf{Discriminator}\\
\midrule
\textbf{Theory 1} & [Prediction in standard exp] & [Test that would fail under T2] \\
\textbf{Theory 2} & [Same observable, alt mechanism] & [Test that would fail under T1] \\
\bottomrule
\end{tabular}
\end{table}

[Move 3, the turn, one sentence; then move 4: the approach walked
through evidence-first, each clause keyed to its section, and the
method given its name here. Model / dataset / canonical configuration
in one sentence.]

[Move 5: contributions.]
\begin{enumerate}[leftmargin=*, itemsep=2pt, topsep=2pt]
\item {[Claim 1: one sentence with the headline number.]}
\item {[Claim 2.]}
\item {[Claim 3.]}
\end{enumerate}

% =============================================================================
\section{Phenomenon}
\label{sec:phenomenon}

[Brief replication of prior work's measurement under our refined
protocol. Include the methodological controls that rule out trivial
confounds. Compare to prior numbers.]

% =============================================================================
\section{Controls: what's actually triggering the effect?}
\label{sec:controls}

This section follows a \emph{discovery arc}: each control tightens the
question for the next experiment.

\subsection{The naive control fails, but for the wrong reason}
[Run a first-pass control. Find the loophole.]

\subsection{The refined control}
[Refine the control to close the loophole. Report the discriminating result.]

\subsection{Per-condition discriminator}
[Generalize to a population of conditions; z-score each against the
refined-control distribution.]

\begin{figure}[t]
\centering
% \includegraphics[width=0.75\linewidth]{figures/fig1_main.pdf}
\caption{Main discriminator figure (placeholder). Regenerated by
\texttt{make\_figures.py}; never hand-edit the exported file.}
\label{fig:main}
\end{figure}

\subsection{Discriminator status}
[Summarize: which theory survives, for what fraction of cases. Explicitly
weaken claims where the data is partial.]

% =============================================================================
\section{Mechanism}
\label{sec:mechanism}

The controls established the trigger is not [theory-1's mechanism]. This
section asks the mechanistic questions that remain.

\paragraph{Q1: [first mechanistic question].}
[One-sentence answer with citation.]

\paragraph{Q2: [second].}
[Answer.]

\paragraph{Q3: [third].}
[Answer.]

\paragraph{Q4: [fourth].}
[Answer.]

\subsection{Synthesis}
\label{sec:synthesis}

[Combined picture: one paragraph summarizing the mechanism story.]

\paragraph{What we explicitly do NOT claim.}
\begin{itemize}[leftmargin=*, itemsep=1pt, topsep=2pt]
\item \emph{Not} ``[over-strong claim 1]'' (falsified by Q$x$'s answer).
\item \emph{Not} ``[over-strong claim 2]'' ([reason]).
\item \emph{Not} ``[over-strong claim 3]'' ([reason]).
\end{itemize}

% =============================================================================
\section{[Behavioral evidence / Cross-task generalization / etc.]}
\label{sec:behavioral}

[Optional § for behavioral evidence beyond the forced-choice measurement.
Include here if relevant to the project.]

% =============================================================================
\section{Related work}
\label{sec:related}

% Prose grouped by theme, not a bullet list; each paper's relation to our work
% stated head-on (the long form lives in doc/related_papers/<citekey>.md).
[Theme 1 in one or two sentences: \citet{example2024} showed X; we differ by Y.]

[Theme 2: \citet{example2025} \ldots]

% =============================================================================
\section{Conclusion}

[1--2 paragraph summary: what we showed, what the field should take
away, what's open.]

% =============================================================================
\FloatBarrier   % flush any remaining body-section figures before references.

\bibliographystyle{plainnat}
\bibliography{references}

% =============================================================================
\appendix
\newpage
\section{Key terms}
\label{app:terms}
\begin{description}[leftmargin=*, itemsep=1pt]
\item[Term 1] [Definition + how it's measured.]
\item[Term 2] [Definition.]
\end{description}

\section{How to read the metrics}
\label{app:metric}
[Primary metric formula + range interpretation + common confusion + z-score formula.]

\section{Per-experiment detail tables}
\label{app:detail}
[Large tables that would clutter §3--§4. Cited as ``→ See Appendix C'' from the body.]

\end{document}
